% !TeX root = ../navier-stokes-manuscript.tex
% ============================================================
% 2.1 Vortex Stretching
% ============================================================

The momentum equation has four pieces:
\[
  \underbrace{\partial_t u}_{\text{local acceleration}}
  +
  \underbrace{(u\cdot\nabla)u}_{\text{nonlinear transport}}
  =
  \underbrace{-\nabla p}_{\text{pressure}}
  +
  \underbrace{\nu\Delta u}_{\text{viscous diffusion}}.
\]
The central regularity problem lies in the competition between nonlinear transport and viscosity.
The hope is simple: viscosity smooths the fluid faster than nonlinear transport can sharpen it.
If that remains true at every scale, the velocity stays smooth.

\subsection{Vortex Stretching}\label{sub:ThePerfectlyStretchingVortex}

Along one material segment of a narrow vortex tube, the surrounding flow pulls in the axial direction. The symmetric part of the velocity gradient records that pull:
\[
  S:=\frac12\bigl(\nabla u+(\nabla u)^{\mathsf T}\bigr),
\]
and if \(e\) points along the segment while \(L(t)\) denotes its length, extension along \(e\) at rate \(\sigma(t)>0\) means
\[
  Se=\sigma(t)e,
  \qquad
  \sigma(t)>0,
  \qquad
  \frac{\dot L(t)}{L(t)}
  =
  e\cdot Se
  =
  \sigma(t).
\]
The same material piece is being drawn out along its own axis.

Because the fluid is incompressible, that extension cannot happen alone. If \(\mathcal V\) is the fixed material volume of the segment, \(A(t):=\mathcal V/L(t)\) its cross-sectional area, and \(r_{\mathrm{eff}}(t):=\sqrt{A(t)/\pi}\), then
\[
  \nabla\cdot u=0,
  \qquad
  \frac{d}{dt}(A L)=0,
  \qquad
  \frac{\dot A}{A}=-\sigma(t),
  \qquad
  \frac{\dot r_{\mathrm{eff}}}{r_{\mathrm{eff}}}
  =
  -\frac{\sigma(t)}{2}.
\]
Length is gained exactly where width is lost. The tube is becoming longer and thinner in one motion.

Its rotation moves inside that same deformation. With \(\omega:=\nabla\times u\), the vorticity obeys along the moving fluid
\[
  \frac{D\omega}{Dt}
  =
  S\omega+\nu\Delta\omega,
  \qquad
  \frac{D}{Dt}
  :=
  \partial_t+u\cdot\nabla,
\]
and when the vorticity aligns with the stretching direction, \(\omega=|\omega|e\),
\[
  S\omega
  =
  \sigma(t)\omega,
  \qquad
  \frac12\frac{D|\omega|^2}{Dt}
  =
  \sigma(t)|\omega|^2
  +
  \nu\,\omega\cdot\Delta\omega.
\]
Now the same positive strain that lengthens the segment also amplifies its aligned rotation, while viscosity enters the same balance without giving a pointwise sign by itself.

That growth sits inside the gradient even while the segment still carries only finite kinetic energy. If the speed stays bounded by \(U_0\) on the fixed-volume material segment \(\Omega_{\mathrm{seg}}(t)\), then
\[
  E_{\mathrm{seg}}(t)
  :=
  \frac12\int_{\Omega_{\mathrm{seg}}(t)}|u(x,t)|^2\,dx
  \leq
  \frac12U_0^2\mathcal V
  <
  \infty,
  \qquad
  \left|\frac12\bigl(\nabla u-(\nabla u)^{\mathsf T}\bigr)\right|_{\mathrm F}
  =
  \frac{|\omega|}{\sqrt2}
  \leq
  |\nabla u|_{\mathrm F}.
\]
So the same bounded-speed segment can keep finite energy while its rotating part strengthens and its radius falls. The gradients rise.

Fix \(t_0\) and let \(\ell:=r_{\mathrm{eff}}(t_0)\). Looking at that width through the \NavierStokes{} scaling produces the comparison family
\[
  u_\lambda(x,t)
  :=
  \lambda u(\lambda x,\lambda^2t),
  \qquad
  p_\lambda(x,t)
  :=
  \lambda^2p(\lambda x,\lambda^2t),
\]
for \(\lambda>1\). At time \(\lambda^{-2}t_0\), the corresponding vortex in \(u_\lambda\) has width \(\lambda^{-1}\ell\). Differentiating the rescaled fields gives
\[
\begin{aligned}
  \partial_tu_\lambda(x,t)
  &=
  \lambda^3(\partial_tu)(\lambda x,\lambda^2t),
  \\
  (u_\lambda\cdot\nabla)u_\lambda(x,t)
  &=
  \lambda^3\bigl((u\cdot\nabla)u\bigr)(\lambda x,\lambda^2t),
  \\
  \nabla p_\lambda(x,t)
  &=
  \lambda^3(\nabla p)(\lambda x,\lambda^2t),
  \\
  \nu\Delta u_\lambda(x,t)
  &=
  \lambda^3\nu(\Delta u)(\lambda x,\lambda^2t).
\end{aligned}
\]
The narrowed vortex has been carried to a finer width, but nonlinear transport, pressure, viscosity, and local acceleration all arrive there with the same \(\lambda^3\) weight. The competition from the opening line has not tilted.

What can still change is the measuring scale. The comparison has altered the width being examined, and a norm can grow or shrink for that reason alone. Let \(\Lambda:=(-\Delta)^{1/2}\). At Sobolev height \(s\), a change of variables gives
\[
  \norm{u_\lambda(t)}_{\dot H^s}^2
  =
  \lambda^{2s-1}
  \norm{u(\lambda^2t)}_{\dot H^s}^2.
\]
Only one height removes that extra factor. When \(s=\tfrac12\),
\[
  H(t)
  :=
  \frac12\norm{u(t)}_{\dot H^{1/2}}^2
  =
  \frac12\norm{\Lambda^{1/2}u(t)}_2^2
\]
is left unchanged by the rescaling itself. The levels on either side still move:
\[
  \norm{u_\lambda(t)}_2^2
  =
  \lambda^{-1}\norm{u(\lambda^2t)}_2^2,
  \qquad
  H_\lambda(t)=H(\lambda^2t),
  \qquad
  \norm{\nabla u_\lambda(t)}_2^2
  =
  \lambda\norm{\nabla u(\lambda^2t)}_2^2.
\]
Kinetic energy drops as the width contracts, while first-derivative energy rises. Finite energy is still compatible with the narrowing segment, and \(H\) is the first level at which any further motion has to come from the equation rather than from changing the ruler.

Apply \(\Lambda^{1/2}\) to the momentum equation and pair with \(\Lambda^{1/2}u\). Write the signed contribution of nonlinear transport as
\[
  \mathcal P_{1/2}(t)
  :=
  -\operatorname{Re}
  \pair
    {\Lambda^{1/2}u(t)}
    {\Lambda^{1/2}\bigl((u(t)\cdot\nabla)u(t)\bigr)}_{L^2}.
\]
The pressure pairing vanishes because \(\nabla\cdot u=0\), and the Laplacian contributes \(-\nu\norm{\Lambda^{3/2}u}_2^2\). The critical-height balance is then
\[
  H'(t)
  +
  \nu\norm{\Lambda^{3/2}u(t)}_2^2
  =
  \mathcal P_{1/2}(t).
\]
Here the original hope has reached the one scale where scaling itself says nothing. If \(H\) rises, nonlinear production has exceeded viscous dissipation there; if \(H\) falls, viscosity has done more.

The same rescaling that took \(\ell\) to \(\lambda^{-1}\ell\) keeps those two contributions matched:
\[
  \mathcal P_{1/2}[u_\lambda](t)
  =
  \lambda^2\mathcal P_{1/2}[u](\lambda^2t),
  \qquad
  \nu\norm{\Lambda^{3/2}u_\lambda(t)}_2^2
  =
  \lambda^2\nu\norm{\Lambda^{3/2}u(\lambda^2t)}_2^2.
\]
The finer comparison still gives neither side a relative gain. What it does give is the contracted time coordinate \(t\mapsto\lambda^{-2}t\), so the same unresolved balance is being read on a shorter clock.

For viscosity, that clock is already present in the linear heat flow. If \(v_t=\nu\Delta v\), then
\[
  \widehat v(\xi,t+\delta t)
  =
  e^{-\nu|\xi|^2\delta t}\widehat v(\xi,t).
\]
A structure concentrated at width \(r\) lives at frequencies \(|\xi|\simeq r^{-1}\), and the multiplier changes by order one when \(\nu|\xi|^2\delta t\simeq1\). The width \(r\) is thus paired with the viscous comparison time
\[
  \tau_\nu(r)
  :=
  \frac{r^2}{\nu}.
\]
This is not the lifetime of the material tube. It is the time diffusion alone needs to alter a structure of width \(r\) by order one. At the finer width,
\[
  \tau_\nu(\lambda^{-1}r)
  =
  \lambda^{-2}\tau_\nu(r),
\]
so the viscous comparison time contracts by exactly the same factor as the full \NavierStokes{} scaling clock. As the width drops, the balance remains level and the time attached to that width drops with it.

Repeat that pairing through the geometric sequence
\[
  r_n:=2^{-n}r_0,
  \qquad
  \tau_n:=\frac{r_n^2}{\nu}.
\]
Then
\[
  \tau_n
  =
  \frac{r_0^2}{\nu}4^{-n},
  \qquad
  \sum_{n=0}^{\infty}\tau_n
  =
  \frac{4r_0^2}{3\nu}
  <
  \infty.
\]
Infinitely many finer width--time pairings fit inside a finite total duration. That by itself is still only a schedule of comparisons. To describe a candidate singularity, one actual solution would have to pass through widths
\[
  r_0>r_1>r_2>\cdots\longrightarrow0
\]
at times
\[
  t_0<t_1<t_2<\cdots\uparrow T_*<\infty,
  \qquad
  t_{n+1}-t_n\asymp\tau_n,
\]
with comparison constants independent of \(n\). The remaining time would then satisfy
\[
  T_*-t_n
  \asymp
  \sum_{k=n}^{\infty}\tau_k
  =
  \frac{4r_n^2}{3\nu}.
\]
By the stage where the core has narrowed to \(r_n\), only order \(r_n^2/\nu\) of comparison time is left. The concentration problem has now become temporal: the same smooth velocity history would have to keep producing the next response on a sequence of shrinking clocks whose total length is finite.
